Aquí sí entra la teoría, porque es la que explica por qué la cota se cumple.

La lectura es una resta que no puede fallar. Solo cuando el peso baja de P hay que recargar, y la recarga sí usa un compare-and-swap, que puede fallar. La clave del paper es qué hacer cuando falla: si otro ya cambió el puntero o ya recargó, me voy; solo reintento si el peso bajó. Como el peso empezó por debajo de P y no puede bajar de cero, como mucho reintento P veces.

La regla de que el taco tenga el doble de vales que hilos garantiza que el peso nunca sea negativo.

El simulador del equipo prueba intercalados al azar y comprueba las reglas en cada paso. El algoritmo del paper no falla nunca; cada variante a la que le quitamos una pieza, sí. Y los intentos con P igual a 4 nunca pasan de 4.
